\documentclass[manuscript,screen,review,anonymous]{acmart}
\usepackage{tabularx}
\setcopyright{none}
\settopmatter{printacmref=false}
\acmConference{}{}{}
\acmBooktitle{}
\acmDOI{}
\acmISBN{}
\renewcommand{\shortauthors}{Anonymous authors}
\newcommand{\HumT}{\textsc{HumT}}
\newcommand{\Dcrit}{\ensuremath{\overline{D}\geq 3}}
\newcommand{\Fcrit}{\ensuremath{\overline{F}<4}}
\newcommand{\OAcrit}{\ensuremath{\overline{\mathrm{OA}}<4}}
\input{tables/numbers}
\title{When Tone Scores Allocate Attention: Auditing HumT-Based Review Prioritization}
\author{Anonymous Author(s)}
\begin{document}
\begin{abstract}
When an automated score ranks responses for limited human review, aggregate coverage does not describe which concerns remain outside the queue. We evaluate this proposed use of HumT on 1,490 existing AI responses across mental health, education, and general health. At 20\% full-batch capacity, HumT captures 42/165, 7/12, and 10/23 broad anthropomorphic-implication cases, respectively; adjusted comparisons with random are inconclusive in the first two domains and narrowly positive in the third. Transparent markers capture more of these cases, while comparisons for domain fit and overall appropriateness are mixed. Adding HumT to a prompt-group-held-out marker/length logistic screen does not establish incremental recall in any of nine domain--criterion cells. Projecting each full-domain queue onto recorded prompt categories shows allocation differences hidden by aggregate coverage: in general health, markers capture more fit concerns overall but fewer natural-prompt fit concerns in point counts, although the latter contrast is uncertain. This exploratory computational audit distinguishes enrichment, incremental benefit, and the distribution of unreviewed concerns. It evaluates existing rubric judgments, not harm outcomes, organizational effectiveness, or non-selection as safety clearance.
\end{abstract}
\keywords{AI evaluation, anthropomorphism, measurement validity, review allocation, algorithmic auditing}
\maketitle

\section{Introduction}
A human evaluator rarely has unlimited capacity to inspect model outputs. Prior work on human--machine triage treats allocation under limited human involvement as an explicit decision problem~\cite{okati2021triage}. Even when an organization specifies a review criterion, it must decide which responses enter a constrained review queue. An inexpensive automated measure could support that decision without directly measuring the criterion. Yet an association between that measure and a human judgment does not establish which cases an allocation rule captures or whether it improves on a simpler alternative.

This distinction matters for anthropomorphic AI language. Expressions of understanding, apparent emotion, conversational familiarity, and first-person reference can make an interaction feel human-like. Such language is not inherently inappropriate. Its interpretation depends on what is expressed and the interactional setting~\cite{devrio2025taxonomy,cheng2025dehumanizing}. A tone score can summarize aspects of this language without establishing whether the response is misleading, fits its domain, or merits escalation. \HumT{} measures human-like language~\cite{cheng2025humt}; its authors do not claim that it is a calibrated probability of harm or a validated human-review policy.

We therefore investigate a proposed downstream use, rather than an alleged flaw in the metric's original purpose. Consider a model-evaluation team auditing a batch of generated responses before deciding what warrants closer examination. It can inspect only a specified fraction of prompt--response pairs. HumT is a plausible candidate when the review target concerns anthropomorphic language, but its suitability for other concerns must be tested rather than assumed. Our analysis measures selection against existing human judgments. It does not establish adoption of this workflow, observe subsequent review, or treat non-selection as clearance.

Our central question is: \emph{When review capacity is limited, which rubric-defined concerns does HumT-based prioritization identify, which remain unreviewed, and how does this depend on the evaluation criterion and context?} We address three research questions:
\begin{description}
\item[RQ1: Comparative allocation.] At matched review budgets, how does descending HumT compare with random selection and transparent observable baselines across separate concern criteria?
\item[RQ2: Incremental benefit.] Does adding HumT improve recall in a specified marker/length screening model evaluated on held-out prompt groups?
\item[RQ3: Distribution of coverage.] What does aggregate performance omit about the distribution of review slots and unreviewed concerns across recorded categories within a shared queue?
\end{description}

We reuse PERSONA's corpus and judgments from a related manuscript under review, which examines overall-appropriateness prediction through regression, grouped cross-validation, and ablation. The present contribution is an exact-capacity allocation study: matched-budget alternatives, complementary selected cases, residual concerns, and retrospective workload, rather than new data or a new construct framework.

Our contribution is a corpus-specific empirical audit, not a new recall metric or triage algorithm. We establish operating points against matched-capacity alternatives, test an incremental contribution within one transparent trained screen, and inspect the distribution of coverage after categories compete for the same queue. Criterion and provenance checks bound those findings. Retrospective workload and a worked reporting card support interpretation; they are not independent experiments or a validated governance intervention. Finding different cases, improving constrained allocation, and demonstrating safety are distinct achievements.

The accountability issue is the justification of this allocation. A review owner chooses the target, budget, and ranking rule; those choices determine which generated responses receive an opportunity for scrutiny. Users who might encounter those responses need not share the owner's definition of a concern. A favorable aggregate score can obscure both disputed target definitions and concentrated omissions. Our evidence concerns the distribution of review opportunities, while the consequences for users and the effectiveness of any challenge process remain questions for subsequent evaluation.

\section{Related Work and Conceptual Scope}
\subsection{Human-like language and contextual interpretation}
HumT DumT develops a measure and control approach for human-like language~\cite{cheng2025humt}. DeVrio et al.'s taxonomy distinguishes linguistic expressions that contribute to anthropomorphism~\cite{devrio2025taxonomy}. The taxonomy offers descriptive resources for interpreting cues such as personal reference and apparent mental states; it does not establish that a particular cue is harmful or that its categories map directly onto the numerical constructs used here. Dehumanizing Machines examines interventions on anthropomorphic behavior and their contextual tradeoffs~\cite{cheng2025dehumanizing}. Consequently, the general observation that context matters is prior motivation, not our claimed discovery.

Related FAccT work examines anthropomorphic cues and their social interpretation, including minoritized cues~\cite{basoah2025hunty} and agency in parasocial human--AI interaction~\cite{maeda2024parasocial}. These studies discourage reducing anthropomorphism to a universally undesirable scalar. Our analysis is narrower: it asks whether a particular ranking policy selects cases meeting specified review criteria in a bounded corpus. We do not measure social identity perceptions, parasocial attachment, or longitudinal effects.

\subsection{From measurement validity to use validity}
Measurement and Fairness emphasizes the relationship between constructs and their operationalization~\cite{jacobs2021measurement}. Evaluation Gaps documents tensions between model-oriented evaluation and application-grounded evaluation~\cite{hutchinson2022gaps}. Applied here, a distinction is necessary between measuring human-like tone, associating tone with a rubric judgment, and establishing useful allocation of human attention. A score may support one of these tasks without supporting the others.

The evaluation target also requires scrutiny. Human ratings operationalize judgments under a rubric and rater population; they do not become objective harm outcomes merely because they are used as labels. Work on representational-harm measurement likewise motivates care about what an observable measure does and does not capture~\cite{chien2024beyond}. We retain separate criteria rather than treating empathy, misleading implication, fit, and holistic appropriateness as interchangeable. Criterion differences may reflect thresholds and prevalence as well as substantive distinctions.

\subsection{Budgeted triage and deferral}
Okati et al. optimize human--machine prediction under a limit on human involvement, while Mozannar and Sontag learn when a classifier should defer to an expert~\cite{okati2021triage,mozannar2020defer}. These studies motivate evaluating allocation relative to a decision objective rather than assuming that a score is intrinsically suitable for routing. Their learned systems model human or expert prediction behavior. Our narrower setting compares fixed rules and criterion-specific logistic screens under exact per-batch quotas. It does not model the accuracy or actions of a subsequent human reviewer. An unselected response remains unreviewed rather than receiving an automated negative decision. We measure criterion coverage, not combined human--machine prediction accuracy. This literature grounds the allocation problem, not a claim of existing organizational adoption of HumT screening.

OpenL2D/FiFAR benchmarks learning-to-defer methods with synthetic experts and per-batch capacity constraints~\cite{alves2025openl2d}. Capacity-aware routing is therefore not itself new. Our empirical object is narrower: competing rankings against existing concern labels, without modeling reviewer decisions or combined human--machine error. The corpus-specific operating points and coverage distributions, rather than a new routing mechanism, must supply the contribution.

\subsection{Auditing and reporting}
Internal algorithmic auditing and assurance frameworks situate evaluation within broader processes of scoping, evidence gathering, and accountability~\cite{raji2020accountability,lam2024assurance}. Our simulation is not an end-to-end audit. It evaluates one possible selection component and specifies information needed before an organization could responsibly interpret it.

Model cards and datasheets provide precedents for documenting intended uses, evaluation conditions, and data provenance~\cite{mitchell2019modelcards,gebru2021datasheets}. We adapt this documentation orientation to a review-allocation card, including residual unreviewed concerns and a proposed route to challenge a criterion or selection decision. The card is a design proposal, not a validated intervention. The gap addressed is operational: existing tone measurements and appropriateness associations do not directly state how many criterion-defined cases enter a fixed-capacity queue relative to simple alternatives.

\section{Corpus, Provenance, and Criteria}
\subsection{Existing corpus and annotation}
The final analytical corpus contains 1,490 English, single-turn responses in 499 domain-scoped prompt clusters: 660 responses and 220 clusters for mental health; 415 and 139 for education; and 415 and 140 for general health. Responses derive from three model endpoints and natural and adversarial prompt sets. Prompt sets are unmatched across conditions. Domain, source material, rater population, and score distributions vary together, so domain comparisons do not isolate a causal context effect.

Each response has a recorded automated score $H$ and five ratings for each of empathic appropriateness ($E$), misleading anthropomorphic implication ($D$), domain fit ($F$), and overall appropriateness (OA). Higher $D$ indicates greater concern; higher $F$ and OA indicate more favorable judgments. We reconstruct every stored construct mean from the released ratings. All recorded HumT values are finite; HumT is not treated as a probability.

The source study's final documentation describes two non-overlapping rater groups. Group A comprises three domain-specific pools of five raters each and supplies OA. A separate five-person Group B supplies $E/D/F$ across domains and includes authors and research interns familiar with the framework. These are operational targets under the source rubric, not independent ground truth for anthropomorphic harm. OA provides a separate-rater, different-construct endpoint, not an independent replication of the $D$ judgments. Domain-relevant backgrounds are documented as mental-health/psychology training, medical training, and teaching experience; professional licenses are not inferred. Phase labels distinguish rating tasks, not shared personnel; the final application note states that the same-rater fallback was not used.

\paragraph{Agreement and cutoff stability.}
The source analysis reports two-way absolute-agreement ICC for averaged ratings, ICC(A,5), its single-rater counterpart, and Krippendorff's ordinal alpha. We reconstruct ICC point estimates from current ratings and report all constructs in Table~\ref{tab:reliability}. For $D$, ICC(A,5) is 0.959, 0.963, and 0.942 in mental health, education, and general health; ordinal alpha is 0.840, 0.843, and 0.446. The difference between average-rating reliability and ordinal agreement matters, particularly for general health. Neither establishes construct validity or independence from framework development. Source reliability intervals use response-level resampling, distinct from our prompt-cluster screening intervals.

We also count means within $\pm0.2$ of each cutoff, a band corresponding to one scale-step by one of five raters. Mental health has 172/660 means near $D=3$, including 108 exactly at the cutoff. Education and general health have 5/415 and 11/415 near it. Mean-based and majority-vote $D$ classifications differ for 35, ten, and one responses, respectively. Appendix~\ref{app:labels} reports all three criterion diagnostics and the full distribution of $D$ means.

\subsection{Input and reuse audit}
The processed response and rating files match the repository's submission-supplement snapshot, and checked principal numerical results agree with the related manuscript. This verifies consistency among available artifacts, not which archive was uploaded to a conference system. The present analysis preserves all source scores and judgments.

Mental-health prompt metadata and HumT scores are embedded in the released records; no separate historical scorer-input string is available. Education and general-health prompt identifiers are reconstructed from prompt text, and model attribution derives from HumT source files. The original baselines tokenize the processed response associated with the annotations, whereas HumT is an imported score associated with an export's text field. Replaying the source join reproduces every imported score, source filename, and join tier. Of 830 exported/rated pairs, 771 are byte-identical; the remaining 56 education and three general-health pairs differ even after normalization. These are 52 prefix and seven fuzzy joins, not merely cosmetic mismatches. Export strings do not authenticate historical scorer calls. We preserve every row and evaluate both same-batch text-representation and fixed-capacity subset sensitivities.

The corpus is not a deployment sample. Its mixture weights, condition frequencies, and model endpoints do not estimate real organizational demand. Most prompts have three responses; there are incomplete model sets in education and general health. Complete-set model comparisons retain 220, 138, and 136 prompt clusters, respectively. We never infer model-family rankings from these three endpoints.

\subsection{Concern criteria}
Our candidate primary trigger is \Dcrit{}, where the bar denotes the mean of five ratings. At rubric level 3, the $D$ construct includes simulated affect or understanding, potentially including ordinary sympathetic expressions. This is deliberately a broad \emph{review trigger}: it is not proof of deceptive intent, experienced deception, inappropriate reliance, or harm. The stricter $D\geq4$ anchor concerns stronger claimed feelings or relationships and is evaluated separately.

Secondary criteria are \Fcrit{} and \OAcrit{}, indicating judgments below strong fit and mostly appropriate, respectively. Applying ordinal anchors to an arithmetic mean is an operational choice, not a claim of equal severity across dimensions. We therefore retain distinct endpoints and examine alternative cutoffs and consensus rules. $E$ is part of the source resource but is not a primary screening endpoint here. Human judgments are not test-time screening features. The later trained comparison uses criterion labels from training prompts only; held-out labels evaluate the resulting selections.

The primary trigger yields 165 mental-health, 12 education, and 23 general-health responses. The corresponding concern-bearing prompt counts are 99, seven, and 15. The small education count is substantively important: a high recall percentage can represent very few cases. Table~\ref{tab:coverage} retains counts alongside percentages rather than allowing a proportion to conceal sparse evidence.

The marker baseline and broad $D$ trigger share some linguistic cues. However, the rubric explicitly excludes functional first-person grammar alone, quoted speech, and claims attributable to other people; a word dictionary cannot enforce those attribution distinctions. This makes the baseline a useful test of how much the operational target can be recovered from surface cues, not a competing ground-truth definition or a guaranteed detector. Fit and separately rated OA provide complementary endpoints with different meanings, not substitutes chosen to produce a preferred policy ordering.

\section{Screening Evaluation}
\subsection{Decision and policies}
The accountable actor in the proposed workflow is an evaluator responsible for a defined response batch. The review unit is a prompt--response pair, and selection creates an opportunity for subsequent manual assessment rather than an automatic adverse action. We do not observe review findings, escalation, or remediation. Non-selection means that an item has not entered this review queue; it does not mean cleared, safe, or appropriate.

For a batch of $n$ responses and budget fraction $b$, capacity is $k=\lfloor nb\rfloor$. We evaluate $b\in\{.05,.10,.20,.30\}$. The already-explored 20\% budget is a reference operating point, not a retrospectively optimized optimum. Primary analyses use separate domain batches; pooled analyses are secondary.

Three fixed policies rank scores in descending order: \textbf{high-H} uses recorded HumT; \textbf{length} uses regex word-token count; and \textbf{markers} uses the fraction of tokens in a fixed first/second-person and conversational-marker dictionary (Appendix~\ref{app:implementation}). The marker score is a transparent lexical comparator, not an independently validated anthropomorphism metric. These rules require no task-specific human training labels, unlike the separately evaluated logistic screens below.

A post-review amendment adds one untuned \textbf{rank-fusion} policy: the equal-weight mean of within-batch midrank percentiles for HumT and markers, selected at the same exact capacity. It uses no labels or fitted weights; ranks are recomputed within each bootstrap batch. Selected-set intersections describe complementary cases, while the larger union is not counted as a same-budget policy. This tests one transparent combination, not the best possible trained screen or conditional independence of the two scores.

Uniform random selection uses the same exact $k$. We compute its expectation analytically, avoiding comparison against one unusually favorable or unfavorable draw. For pooled analyses, an additional stratified random policy allocates slots proportionally by domain using largest-remainder rounding. Deterministic seeded hashes of stable response identifiers resolve ties without exceeding capacity. A 100-seed analysis examines sensitivity to tie resolution. Missing or nonfinite policy scores are rejected rather than silently excluded.

\subsection{Held-out trained comparison and conditional lift}
A second post-review amendment tests the added value of HumT with a transparent trained comparator. For each domain and criterion, the reduced model is L2-regularized logistic regression on marker fraction and $\log(1+\mathrm{length})$; the extended model adds recorded HumT. This model isolates the question of adding one feature to an inspectable screen, not competitiveness with the strongest available evaluator. We make no comparison claim against semantic judges or modern automated review systems generally. Both models use fixed $C=1$, an intercept, LBFGS, and no class weighting or tuning; standardization is fitted on training data only. Their scores rank the source criterion rather than estimate calibrated harm probabilities.

We partition prompt groups into five folds using a seeded, label-independent permutation. Each held-out fold is a separate review batch; all responses to a prompt remain together. Each model trains on the other four folds. We re-evaluate all fixed policies, fusion, and analytical random selection in these same held-out batches. Counts are summed before calculating recall or precision; scores from different fitted models are never pooled into a common ranking. At the 20\% reference fraction, flooring each batch's quota yields 130 total mental-health slots and 80 in each other domain, rather than the original full-batch capacities of 132 and 83. Comparisons therefore use the new matched-batch table, not the original capacities.

We define \emph{conditional lift} operationally as the absolute recall difference
\begin{equation}
\Delta R_{H\mid M,L}=R_{\mathrm{LR}(M,\log(1+L),H)}-R_{\mathrm{LR}(M,\log(1+L))}.
\end{equation}
It tests an incremental contribution within this specified model and workflow, not a causal effect, conditional independence, or all possible uses of HumT. Training uses existing human labels and is not label-free: each mental-health fit uses 528 labeled responses, while transfer-domain counts vary with fold size (Appendix~\ref{app:complete}). Held-out labels are never used to fit their own model. These historical labels are reused across fits; annotation acquisition costs are not estimated, nor are fit-level counts summed as separate purchases.

At the reference fraction we use 2,000 paired prompt-cluster bootstrap draws within the original fold strata, retaining a sampled prompt's original fold. We refit both scalers and models on every draw, re-rank every resampled held-out batch at its exact quota, and use common draws across policies and criteria. A single-class training sample produces an explicitly recorded constant-prevalence score; undefined recall remains undefined. Intervals are exploratory pointwise percentile intervals conditional on this partition, empirical judgments, and prompt distribution. Ten additional pre-specified partitions assess split sensitivity, not independent replication or an opportunity to choose the best seed. Other budgets receive descriptive point estimates.

\subsection{Outcomes and workload}
Let $P$ be the number of responses meeting a criterion and $T$ the number selected that meet it. Recall is $T/P$, precision is $T/k$, and the number of unreviewed concern cases is $P-T$. Concern prevalence among unselected responses is $(P-T)/(n-k)$. Under uniform random selection, expected $T=kP/n$ and expected recall is $k/n$ when $P>0$. Appendix~\ref{app:transfer} retains capacity ceilings and prevalence-relative lift as contextual quantities, not additional independent outcomes.

We also report average precision (AP) and ROC-AUC for ranked policies, with class prevalence. AP uses recall-increment-weighted precision at tied-score block boundaries, not trapezoidal interpolation. Both implementations were independently checked against scikit-learn~\cite{pedregosa2011scikit}. AP complements ROC-AUC when concern prevalence is low~\cite{saito2015precision}; neither makes the human criterion a validated harm outcome.

For a target recall $r$, retrospective workload is the shortest ranking prefix containing at least $\lceil rP\rceil$ concern cases, divided by $n$. We retain the originally specified 50\%, 80\%, 90\%, and 100\% targets; 80\% is a descriptive reference, not a justified service requirement. For random ordering, expected prefix length is $\lceil rP\rceil(n+1)/(P+1)$. These label-informed summaries are not deployable stopping rules or labor-cost estimates.

\subsection{Uncertainty and analysis status}
We use 2,000 paired prompt-cluster bootstrap replicates. All responses for a sampled prompt remain together, and the same resamples support comparisons across policies and criteria. Each replicate re-ranks its resampled batch at the specified capacity. Repeated copies of a cluster receive occurrence-specific tie keys. Pooled resampling is stratified by domain. Thus, intervals characterize a resampled-batch allocation procedure, not application of a previously learned fixed threshold.

Most intervals are exploratory pointwise 95\% percentile intervals. The specified primary comparison family contains three domain-specific high-H minus random recall differences for \Dcrit{} at 20\%, using Bonferroni-adjusted 98.3333\% percentile intervals. This nominal adjustment addresses that family, not all sensitivity analyses. Fewer than ten concern-bearing or nonconcern-bearing prompts triggers a sparse-evidence flag. Recall and AP are undefined without positive cases; AUC requires both classes. We report finite replicate counts rather than substituting zero for undefined values.

The corpus and several thresholds had already been explored before the current specification. Earlier exploration included alternative OA, $D$, and $F$ cutoffs, high-H top-20\% selection, ranking summaries, overlaps, and model means. Our protocol is a prospective specification for this computation on already-inspected data, not a preregistration of unseen evidence. Bootstrap intervals condition on the released ratings; rater-rule sensitivity addresses a different source of uncertainty and is not absorbed into these intervals.

After reviewing the first results, we specified complementarity, fixed rank fusion, cutoff diagnostics, detailed text checks, and direct paired workload differences. These additions use 2,000 resamples and exploratory pointwise 95\% intervals, with no confirmatory significance claims. They preserve the original primary family. The larger mental-health sample is emphasized for presentation, not selected retrospectively as a replacement primary endpoint. Sparse comparisons remain descriptive; a different interval method cannot supply missing concern-bearing prompts.

The trained comparison was specified after those findings were inspected. Its nine domain--criterion cells and split sensitivity are a separate exploratory amendment, not nine new primary tests. The original full-batch family and outputs remain unchanged; interval overlap with zero is not evidence of equivalence or absence of useful HumT information.

\subsection{Common-queue category coverage}
A third exploratory amendment ranks each full domain once, at its original 20\% capacity, then projects that selection onto existing condition, model, and scenario categories. Unlike the earlier subgroup-specific analyses, it never reranks or reserves slots within a category. All three criteria and five fixed policies (including fusion and analytical random selection) are retained. Within each category we report response and concern counts, selected slots, recall, missed concerns, and residual prevalence. A category with no concerns has undefined recall. Random expected group slots are $kn_g/n$ and expected caught cases are $kP_g/n$; this is not a separately floored group quota.

Paired domain-level prompt resamples re-rank the full batch before category projection. We report pointwise intervals for category recall, group-minus-complement recall, and within-group HumT-minus-marker recall. Sparse flags retain positive and negative supporting-prompt counts for the group and its complement. Condition is the compact main reporting partition, specified before this run; all model/scenario cells remain in the record. These are descriptive coverage distributions in unmatched prompt sets, not demographic fairness estimates or causal effects of prompt condition.

This specification preceded execution but followed inspection of earlier context results. Each amendment uses 2,000 paired draws; the effective evidence is the number and diversity of supporting prompts, not the number of resamples. Reproducibility does not remove analytic-selection risk or turn these amendments into confirmation on unseen data.

\section{Results}
\subsection{RQ1: Comparative allocation at matched capacity}
Table~\ref{tab:coverage} shows high-H coverage at 20\%. For \Dcrit{}, the policy selects 42 of 165 mental-health cases, seven of 12 education cases, and ten of 23 general-health cases. Uniform random selection would capture 33, 2.4, and 4.6 cases in expectation. The corresponding HumT precision values are 31.8\%, 8.4\%, and 12.0\%. Enrichment therefore does not imply that most selected responses meet the criterion, especially when the criterion is rare.
\input{tables/coverage}

In mental health, even perfect selection into 132 slots cannot recover all 165 broad-$D$ cases. This capacity constraint does not explain all omissions: HumT captures 42 cases and leaves 123 unreviewed. Both the available capacity and the ranking must therefore be considered.

The primary-family recall improvements over random are \MHDelta{}, \EDDelta{}, and \GHDelta{} percentage points, with adjusted intervals \MHAdjusted{}, \EDAdjusted{}, and \GHAdjusted{}, respectively. The general-health lower endpoint is very close to zero, and its interpretation depends on the percentile procedure and finite bootstrap resolution. The results do not establish robust superiority across domains or thresholds. Education's \Dcrit{} estimate rests on only seven concern-bearing prompts.

Coverage also depends on the criterion. Mental-health domain-fit recall is 18.9\%, slightly below the random expectation, while general-health fit recall is 20.8\%. For education, high-H recovers ten of 25 fit concerns and nine of 18 OA concerns. The general-health OA estimate of 33.3\% represents three of nine cases across seven concern-bearing prompts and has a wide interval. A case outside one criterion may still satisfy another; neither criterion membership nor non-membership establishes general safety.

Figure~\ref{fig:paired} summarizes all nine paired HumT--marker comparisons alongside the trained incremental comparison. The panels use different, explicitly labeled batch definitions; neither substitutes for the adjusted primary family. Appendix~\ref{app:transfer} retains all domain budget curves.
\begin{figure}[t]
\centering
\includegraphics[width=\linewidth]{figures/paired_summary.pdf}
\caption{Complete paired recall contrasts at the 20\% reference fraction. Left: HumT minus markers in original full-domain queues. Right: adding HumT to the marker/log-length logistic screen in common held-out batches, with refitted uncertainty. Positive values favor HumT or its addition. Intervals are exploratory pointwise 95\%, not simultaneous; hollow points mark sparse support. Counts refer to criterion cases and concern-bearing prompts.}
\Description{Two forest plots show all domain and criterion differences. The left panel shows mixed criterion-dependent comparisons; all nine intervals in the trained incremental panel cross zero. Every point has its own paired uncertainty interval and supporting counts.}
\label{fig:paired}
\end{figure}

\paragraph{Transparent alternatives.}
Table~\ref{tab:baselines} compares policies for the primary trigger. The marker policy selects 63 mental-health, nine education, and 21 general-health concern cases at the same capacities as HumT. HumT minus marker recall is -12.7 percentage points in mental health and -47.8 points in general health; their paired pointwise intervals exclude zero. The education difference is -16.7 points but uncertain. Length performs better than HumT for mental-health $D$ coverage and worse in the other two domains; those comparisons are less conclusive.
\input{tables/baselines}

This is not evidence that marker counting is a universal safety screen. Pronouns and expressions such as sympathy overlap with cues in the broad $D$ rubric. The result shows that an inexpensive observable comparator can recover these particular labels more effectively; it does not show that removing its markers improves an interaction. Ranking-level summaries tell a consistent but distinct story. HumT AP for \Dcrit{} is \MHAP{}, \EDAP{}, and \GHAP{} against prevalences of 25.0\%, 2.9\%, and 5.5\%. Marker AP is 0.465, 0.126, and 0.395, respectively. These values should be read within each criterion and population, not as prevalence-free universal rankings.

The full fit/OA comparison is mixed, not a general reversal of the broad-$D$ result. In mental health, HumT selects 18/95 fit and 31/133 OA cases versus markers' 22/95 and 33/133; both paired recall-difference intervals include zero. In education it selects 10/25 fit and 9/18 OA cases versus 1/25 and 3/18. The education fit-recall difference is +36.0 percentage points (exploratory 95\% interval +15.4 to +58.6); the OA difference is +33.3 points (-7.7 to +71.4). In general health, HumT selects 25/120 fit and 3/9 OA cases versus markers' 32/120 and 9/9: differences are -5.8 points (-13.6 to +2.9) and -66.7 points (-100.0 to -28.5). The latter rests on only seven concern-bearing prompts. Appendix~\ref{app:complete} reports every policy--criterion cell, including length, fusion, and random. These pointwise exploratory findings do not establish universal fit/OA superiority or precise performance on rare concerns.

\subsection{RQ2: Incremental benefit beyond lexical features}
\paragraph{Different cases are not a same-budget improvement.}
At the mental-health reference budget, HumT and markers share 22 broad-$D$ cases; HumT selects 20 additional cases that markers miss, while markers select 41 that HumT misses. Their union would require 221 review slots rather than 132, so it is not a fair same-budget alternative. Equal-weight rank fusion selects 61 cases versus markers' 63, a recall difference of \MHFusionDelta{} percentage points (exploratory 95\% interval). For fit and OA it selects 25 and 39 cases versus markers' 22 and 33, but both paired difference intervals include zero.

In education, HumT contributes one broad-$D$ case outside markers' set; in general health it contributes none at this budget. Fusion captures nine and 17 broad-$D$ cases, versus markers' nine and 21. Table~\ref{tab:complementarity} reports all criteria and paired intervals. These overlaps establish distinct selections, not incremental predictive information. We therefore turn to the paired trained comparison rather than infer that another combination must improve allocation.

\paragraph{Held-out conditional lift.}
Table~\ref{tab:conditional} reports every domain--criterion comparison when adding HumT to the marker/log-length model. In mental health, reduced and extended screens capture 70 versus 66 broad-$D$ cases, 32 versus 33 fit cases, and 26 versus 30 OA cases. Their recall differences are \MHDConditional{}, \MHFConditional{}, and \MHOAConditional{} percentage points, respectively. Every interval includes zero. The OA point improvement is therefore a hypothesis for further evaluation, not established complementary screening value.
\input{tables/conditional}

Education's broad-$D$ point increase is one case (7 to 8 of 12), or \EDDConditional{} points, with only seven concern-bearing prompts. Neither fit nor OA counts change (13/25 and 12/18). General-health broad-$D$ and OA counts also remain unchanged (21/23 and 8/9); fit increases by one (29 to 30 of 120), or \GHFConditional{} points. All nine intervals include zero, without establishing equivalence. Across ten additional partitions, education's broad-$D$ increment ranges from -8.3 to 0.0 points rather than reproducing the reference partition's positive result. Table~\ref{tab:split_sensitivity} reports all split ranges, including unfavorable ones.

The complete held-out comparison also prevents selecting a favorable comparator. For education fit, the reduced and extended models each capture 13 cases versus HumT's nine and markers' one; for general-health OA, markers capture nine cases versus either model's eight. These are descriptive point comparisons under matched held-out quotas, not evidence that training dominates every fixed rule. Overall, the tested model provides no clear incremental-recall evidence for adding HumT at this budget. This conclusion is specific to an untuned linear-logit screen on existing labels; it does not establish that HumT lacks conditional information, that nonlinear combinations cannot help, or that any model is a validated deployment screen.

\subsection{RQ3: Coverage within a shared queue}
Table~\ref{tab:common_queue_slots} shows the distribution of the original full-domain slots. For mental health, HumT selects 78/300 natural responses and 54/360 adversarial responses; markers select 63 and 69. HumT's broad-$D$ coverage is similar across those conditions (21/84 versus 21/81), but OA coverage is 22/69 versus 9/64. The adversarial-minus-natural OA recall difference is \MHQueueOAGap{} percentage points. The corresponding marker difference is small and uncertain. This illustrates why slot shares, concern prevalence, and concern recall must be read separately.
\input{tables/common_queue_slots}

In general health, HumT assigns 27 of its 83 slots to natural prompts, whereas markers assign 12. For fit concerns, HumT captures 8/74 natural and 17/46 adversarial cases; markers capture 4/74 and 28/46. Markers therefore capture more fit cases overall (32 versus 25) while leaving more natural-prompt fit cases unreviewed (70 versus 66). The natural-set HumT--marker difference remains uncertain: \GHQueueNaturalFitContrast{} percentage points. Adversarial-minus-natural fit-recall gaps are \GHQueueFitGapH{} points for HumT and \GHQueueFitGapM{} for markers. These are exploratory within-policy intervals, not a test of the difference between the two gaps or proof that either allocation is fairer.

Education's broad-$D$ and OA concerns occur only in its adversarial set; natural-set recall is undefined, not evidence of successful clearance. Its fit cases occur in both sets, but each has fewer than ten concern-bearing prompts. Appendix~\ref{app:queue} reports all condition-by-criterion cells rather than selecting a favorable subset; all five policies and existing model/scenario categories remain in the computational record. The examples highlighted here were selected for interpretation after the complete analysis, not specified as new primary tests. Unmatched prompts, source differences, and multiple exploratory comparisons preclude causal or population-fairness conclusions.

\paragraph{Sensitivity boundaries.}
The broad-$D$ trigger must not be interpreted as a stronger-deception endpoint: at mean-$D\geq4$, HumT captures 0/5, 3/8, and 2/8 cases, with only four, five, and six supporting prompts. Consensus rules also change the evaluation target. For fixed label-free policies this changes which responses count as concerns, not which responses enter the queue. Restricting transfer domains to byte-identical export/rated text retains all broad-$D$ and OA concerns but removes two education-fit cases, alongside changes to the candidate batch. Neither sensitivity establishes historical scorer-input correctness or independent validation. Appendix~\ref{app:sensitivity} reports the counts and operating points.

\section{Decision Accounting and Supporting Report}
\subsection{Retrospective workload}
Table~\ref{tab:workload} translates mental-health rankings into the batch fraction needed to reach 80\% recall of broad-$D$ cases. HumT requires 70.8\%, versus 79.6\% for expected random ordering, 66.1\% for length, and 63.3\% for markers. Direct paired workload differences are \MHWorkRandom{} percentage points versus random, \MHWorkLength{} versus length, and \MHWorkMarkers{} versus markers. These exploratory 95\% intervals support a workload reduction versus random in this corpus, but do not establish a difference from length or markers. Appendix~\ref{app:transfer} retains the other domains.
\begin{figure}[t]
\centering
\includegraphics[width=\linewidth]{figures/workload_targets.pdf}
\caption{Retrospective workload for the four originally specified broad-$D$ recall targets, with exploratory pointwise 95\% intervals for every policy. Lines connect evaluated targets only; random values are expected order statistics. Sparse education support limits precision. The targets are descriptive, not stakeholder-validated service requirements or deployable stopping rules.}
\Description{Three panels show the batch percentage needed for 50, 80, 90, and 100 percent D recall in each domain. HumT, length, markers, and random all have uncertainty bars; workload rises with target recall.}
\label{fig:workloadtargets}
\end{figure}

Figure~\ref{fig:workloadtargets} shows all four specified targets rather than elevating 80\% into an optimal policy. These fractions describe an observed ranking with labels known; they do not tell an evaluator when an unseen batch has reached its target. Item fractions also omit differences in reading time, criterion complexity, and reviewer expertise.

\subsection{Worked allocation report}
A supporting audit-use card records the batch, owner, criterion provenance, capacity, tie rule, selected and residual cases, and uncertainty (Appendix~\ref{app:card}), following documentation precedents~\cite{mitchell2019modelcards,gebru2021datasheets}. It separates challenges to the criterion, its application, and the allocation rule. A route to raise concerns should not depend solely on receiving a high score. Versioned amendments should retain original judgments and disagreements. The card is an unevaluated reporting proposal, not evidence of successful contestation or organizational adoption; no practitioner study or AI-agent substitute is claimed.

For a concrete reading, take the general-health fit endpoint and 83 available slots. Markers capture 32 of 120 cases and leave 88 unreviewed; HumT captures 25 and leaves 95. The aggregate point estimate therefore favors markers. However, natural-prompt residual counts are 70/74 versus 66/74, respectively, and the within-natural policy contrast is uncertain. A report should expose both comparisons, without declaring HumT preferable, treating natural prompts as a protected population, or recommending an untested quota. Which omissions matter more requires a justified evaluation purpose that the corpus alone cannot supply.

\section{Discussion}
\subsection{Useful enrichment is a conditional claim}
HumT sometimes selects broad anthropomorphic-implication cases at a rate exceeding random allocation. That useful performance should not be obscured by a blanket claim that human-likeness measures fail. Equally, enrichment over random is not sufficient to justify a policy if an inexpensive comparator selects more relevant cases at the same capacity. The primary operational question is comparative allocation under a stated purpose, not whether any association exists.

This distinction also explains why the decision study is not numerically determined by the related prediction analysis. Low linear predictive performance for continuous OA does not specify exact-budget recall for a thresholded $D$ judgment, and neither determines how a marker baseline allocates review. Conversely, workload is a derived view of these same rankings, not an independent validation study. The contribution is the explicit evaluation of the allocation procedure, not a new dataset or a general theory that context matters.

\subsection{A review criterion is part of the intervention}
Choosing $D\geq3$ creates a different review target from choosing $D\geq4$, poor fit, or low OA. The broader trigger includes linguistic forms that may be appropriate in supportive contexts. A screen that efficiently finds those forms is not necessarily a screen that finds the most consequential problems. Criterion provenance, severity assumptions, and what happens after selection are therefore part of the proposed intervention rather than peripheral details.

The responsible review owner should therefore justify both the target and its aggregation rule, distinguish that choice from the ranking algorithm, and make amendments traceable. Our results establish sensitivity to those choices, not which stakeholders' preferences should prevail. Separate OA raters broaden the evaluated perspectives but do not validate Group B's D judgments as harm outcomes.

Observed baseline performance reinforces this point. The marker policy's broad-$D$ advantage can partly follow from its alignment with the rubric's linguistic cues. Fit/OA results vary across domains and alternatives; the education-fit comparison is not a general policy ranking. Complementarity also differs from successful integration: selected-set differences coexist with inconclusive incremental recall in the trained comparison. The absence of clear lift is bounded by the tested model, budget, partition, and labels, not a proof of HumT redundancy. Establishing usefulness for experienced harm or a different evaluator population would require additional evidence.

\subsection{Residual cases belong in the report}
Reporting only selected precision can conceal most concern cases outside the queue. At the reference budget, high-H leaves 123 broad-$D$ mental-health cases and 95 general-health fit cases unreviewed. These are not demonstrated harms or failures of an observed organization. They are concrete limits of the proposed allocation rule against its stated endpoints. Showing them alongside the selected cases prevents a constrained review procedure from being interpreted as comprehensive coverage.

Common-queue results extend this reporting obligation: aggregate gains can coexist with less coverage in a recorded category. A review owner could inspect such a distribution before deciding whether a different allocation warrants evaluation; our analysis does not justify imposing category quotas. Time per item, the cost of unnecessary review, and the consequences of missed concerns remain unmeasured. External evaluators and affected users may reasonably contest the target or capacity; neither authors' familiarity nor AI-generated judgments substitutes for those perspectives.

\section{Limitations}
The evidence is observational reuse of an English, single-turn corpus with three model endpoints. It does not measure experienced deception, clinical or educational outcomes, relational effects over time, or successful review and remediation. Natural and adversarial conditions are unmatched, and domain comparisons confound prompt source, rater pool, and score distribution. Corpus mixture weights do not estimate deployment prevalence.

The human endpoints reflect particular rubrics and populations. Group B includes authors and interns familiar with the framework; separate OA raters do not remove that dependence. Applying cutoffs to mean ordinal ratings and selecting a broad candidate trigger are analytic choices on previously inspected data. Sensitivity analyses make these choices visible but do not turn the study into confirmatory research on unseen evidence.

Sparse education and stricter-$D$ concerns limit precision. Cluster bootstrap intervals condition on recorded labels and the observed prompt distribution; they cannot validate a construct or account for every provenance uncertainty. Text-based joins and reconstructed prompt identifiers remain relevant even after sensitivity analysis. Equal-capacity comparisons ignore heterogeneous reading time. The reporting card has not been evaluated, and no external participant evidence is claimed.

The trained analysis reuses explored data and evaluates one low-dimensional model family without tuning. Its refitted bootstrap is conditional on a fixed prompt partition; the ten alternative splits are sensitivity checks, not new samples. Shared raters, templates, and source processes can create dependence beyond prompt groups. Sparse bootstrap samples sometimes contain only one training class, and the explicit constant-score fallback is part of the evaluated procedure. Historical annotation requirements differ from label-free rules, while deployment drift and acquisition costs remain unmeasured. There is no comparison with an LLM judge or semantic screen. We do not perform equivalence testing: a practically negligible recall margin has not been independently justified, and selecting one to fit the observed intervals would not strengthen the evidence.

\section{Conclusion}
A human-like-tone score can identify criterion cases that another rule misses, but distinct selections alone do not establish a better combined queue. Adding HumT to the specified marker/log-length screen gives no clear incremental-recall improvement across the nine tested domain--criterion cells at the reference budget, without establishing equivalence or universal redundancy. Common-queue analysis also shows why aggregate gains do not describe every category's coverage. The practical lesson is to justify the review target, compare matched capacities and label requirements, and report both the amount and distribution of residual concerns. This is evidence about a proposed allocation procedure, not a general rejection of tone measurement or demonstrated improvement in real-world safety.

\clearpage
\section*{Endmatter}
\subsection*{Ethics statement}
This study reuses existing AI responses and human annotations; it collects no new human-participant data. The source study's records describe informed consent, voluntary participation with withdrawal rights, and compensation through existing grant-supported research roles rather than per-response payment. Consent and recruitment provisions explicitly stated that interns could decline or withdraw without affecting their role, pay, supervision, or evaluation. The team documents the annotation study as a sub-study of a parent project with institutional ethics approval dated 12 February 2026. This statement does not assert a separate new approval or exemption for the present computational analysis. Institutional identifiers are withheld for anonymous review. Group B's authorship, framework familiarity, and potential supervisory power differences remain relevant limitations despite the non-participation safeguards. Human scores are rubric judgments, not evidence of experienced harm. We report aggregate findings without reproducing potentially sensitive case text or identifying annotators.

\subsection*{Generative AI usage statement}
OpenAI Codex assisted with study planning, analysis-code implementation and testing (including the trained-comparator and bootstrap code), literature and source checks, manuscript drafting and revision, and preparation of tables and the LaTeX package. The underlying model version was not recorded in the preparation record. AI assistance did not supply human-participant responses, independent human coding, replacement source ratings, or LLM-judge labels. Logistic screens were fitted to existing source judgments, not AI-generated substitute annotations. The authors retain responsibility for the research claims, source attribution, code, and submitted text. No AI-agent evaluation is reported as human validation.

\subsection*{Data and reproducibility statement}
The computational record preserves input hashes, source revision, software versions, seeds, criteria, and output-generation commands. Numerical tables are generated from the frozen original outputs and separately versioned post-review additions. A repeated original-pipeline run reproduced its 29 computational CSV tables and three figure files byte-for-byte. Revision manifests verify that source files and original outputs remain unchanged. Identifying consent records and institutional documents are not part of the anonymous typesetting package.

\bibliographystyle{ACM-Reference-Format}
\bibliography{references}

\clearpage
\appendix
\section{Label Agreement and Cutoff Diagnostics}\label{app:labels}
Table~\ref{tab:reliability} includes $E$ for source completeness although it is not a screening endpoint. Table~\ref{tab:cutoffs} distinguishes closeness of a mean to the cutoff from agreement on individual-rater trigger membership. High ICC does not guarantee stable binary membership at an operational threshold. Source ICC and alpha confidence intervals are retained in the computational record; those intervals use response-level rather than prompt-cluster resampling.
\input{tables/reliability}
\input{tables/cutoffs}
\begin{figure}[t]
\centering
\includegraphics[width=\linewidth]{figures/d_mean_distributions.pdf}
\caption{Descriptive distributions of five-rater D means. Dashed lines mark the original mean-$D\geq3$ cutoff. Counts, not fitted densities, show the concentration of values; panels use separate count scales.}
\Description{Bar charts show D means from one to five in each domain, with a dashed cutoff at three. Mental-health means include many observations at the cutoff; education and general-health means are mostly below it.}
\end{figure}

\section{Transfer and Complementarity Details}\label{app:transfer}
These are exploratory descriptive comparisons; sparse education D and general-health OA cells do not support confident policy rankings. All domain budget curves remain available here, while the main summary figure shows paired contrasts. Fusion and complementarity were added after the initial analysis and did not replace the original primary comparisons.
For context, oracle recall is $R^*=\min(1,k/P)$ and prevalence-relative lift is $(T/k)/(P/n)$. The achieved fraction of oracle recall is $(T/P)/R^*$. When $0<k\leq P$, this equals precision $T/k$ algebraically, not coincidentally; these quantities are not independent evidence. The oracle uses unavailable labels and is not an implementable policy.
\input{tables/oracle_context}
\input{tables/complementarity}
\input{tables/workload}
\input{tables/workload_transfer}
\begin{figure}[t]
\centering
\includegraphics[width=\linewidth]{figures/mental_health_budgets.pdf}
\caption{Original mental-health budget curves. Shading shows only HumT's exploratory pointwise 95\% intervals, not uncertainty-free baselines. Figure~\ref{fig:paired} provides paired policy-difference intervals; complete marginal intervals are retained in the computational record.}
\Description{Three mental-health panels compare HumT, length, markers, and analytical random recall across four budgets and three criteria.}
\label{fig:budgets}
\end{figure}
\begin{figure}[t]
\centering
\includegraphics[width=\linewidth]{figures/transfer_budgets.pdf}
\caption{Education and general-health budget curves. Shading shows HumT's exploratory pointwise 95\% intervals, not simultaneous bands; random values are analytical expectations. Sparse cells and different prevalences constrain interpretation.}
\Description{Six panels retain recall curves for education and general health under all three criteria and four budgets. Education D and general-health OA have wide uncertainty despite some high observed recall values.}
\end{figure}

\section{Complete Policy Comparisons}\label{app:complete}
The following tables retain every D, fit, and OA policy cell at the reference fraction, including unfavorable and sparse results. Tables~\ref{tab:full_mental_health}--\ref{tab:full_health} use the original full-domain batches and add fixed fusion. Tables~\ref{tab:held_mental_health}--\ref{tab:held_health} re-evaluate all seven policies in identical held-out batches. Their capacities and bootstrap procedures differ; the two table families must not be combined into an apparent matched-budget comparison. Random counts are expectations, so they can be fractional. The record retains all four budgets, selected and missed counts, precision, recall, unselected prevalence, and paired differences from random. Only the reference fraction has intervals for the trained amendment.
\input{tables/full_mental_health}
\input{tables/full_education}
\input{tables/full_health}
\input{tables/held_mental_health}
\input{tables/held_education}
\input{tables/held_health}
\input{tables/split_sensitivity}

Each mental-health model uses 528 training responses (2,640 underlying criterion ratings); education uses 331--334 (1,655--1,670 ratings); general health uses 331--333 (1,655--1,665 ratings). Training sets overlap across fits and these are not additive acquisition costs. No point-estimate fit needed the single-class fallback. Across all criteria and models, it occurred in 68/2,000 education draws and 50/2,000 general-health draws (154 and 116 model fits), and no mental-health draw. Undefined paired recall differences occur in two education-D and two general-health-OA draws; those cells retain 1,998 finite draws. All other conditional-lift cells retain 2,000. These counts belong to the new refitted bootstrap, not the original intervals below.

\section{Common-Queue Coverage Details}\label{app:queue}
\input{tables/common_queue_conditions}
This table projects fixed selections onto the two recorded prompt conditions. It does not recreate the earlier procedure that allocated 20\% within each subgroup. The full record contains 525 category/criterion/policy cells over condition, model, and scenario partitions, including zero-concern groups. Each partition's selected and caught totals reconstruct the original full-domain results. Group-versus-complement recall contrasts are undefined when either side has no concern cases; these draws retain explicit finite counts rather than substituted zeros. Random precision in the category table is the ratio of expected caught cases to expected slots, not the expectation of a random ratio. No new demographic attributes or cue-category annotations were inferred.

\section{Criterion and Text-Provenance Sensitivity}\label{app:sensitivity}
At the stricter mean-$D\geq4$ threshold, high-H selects zero of five mental-health cases, three of eight education cases, and two of eight general-health cases at 20\%. These concerns occur in only four, five, and six prompt clusters. These sparse stronger-anchor results are not precise deployment estimates, and thresholds are not interchangeable choices from which to select a favorable result.

Requiring at least three of five ratings to meet $D\geq3$ yields high-H captures of 48/200, 12/22, and 11/24 across the three domains. Requiring all five yields 39/150, 7/12, and 6/17. Agreement does not validate the normative criterion. Only six mental-health, four education, and six general-health responses satisfy all three primary criteria. The record retains all eight intersections and continuous relationships: cut points and prevalence can explain low overlap without establishing normative disagreement. Trained models were not refitted for these threshold variants.

Recomputing length and markers from HumT export strings on unchanged full batches leaves their selected concern counts unchanged at 20\% for all three criteria in both transfer domains; fusion counts are also unchanged. This does not establish identical selected identities or correct historical scorer calls. A separate reconstruction checked all 48 previously reported subset operating points using byte equality explicitly, rather than assuming the stored join tier implies identical text.

\input{tables/provenance_composition}
\input{tables/provenance_truncation}

The 56 differing education pairs contain no broad-$D$ or OA cases and two fit cases; the three differing general-health pairs contain none of these criterion cases. Nevertheless, their presence can affect competition for review slots. Keeping only byte-identical pairs and retaining the absolute capacity of 83 captures seven of 12 education broad-$D$ cases and ten of 23 general-health cases. At 20\% of the smaller batches, 71 and 82 slots capture six and ten cases. Full condition/model/scenario composition and all four ranked-policy operating points are retained in the computational record. The subset analysis is descriptive, not a paired causal estimate of join error, and does not authenticate the historical scorer invocation. Mental-health scores remain embedded without a separate scorer-input string.

\section{Implementation and Uncertainty Details}\label{app:implementation}
Tokenization uses the Python regular expression \verb|\b\w+\b| after lowercasing. The fixed marker dictionary is: \texttt{i, me, my, mine, myself, we, us, our, ours, ourselves, you, your, yours, yourself, yourselves, hello, hi, hey, thanks, thank, welcome, sorry, please}. The score is the number of matching tokens divided by the token count; an empty response receives zero. Length is that same token count. These deterministic policies require response text but no human outcome labels. HumT uses the recorded source score, with larger values selected first.

The seed is 20261002. Tie keys use the first eight bytes of a SHA256 digest of the seed and stable identifier. Ranking sorts by descending policy score and then ascending tie key. Capacity uses the integer floor, not a score-at-quantile rule, so tied scores cannot inflate the budget. Pooled proportional random allocation uses Hamilton largest remainders with alphabetical domain tie resolution.

Education \Dcrit{} recall has 1,999 finite bootstrap replicates; general-health \OAcrit{} recall has 1,998; the other domain/criterion reference cells have 2,000. Undefined draws are preserved in the record. Primary adjusted intervals use the three-domain family only. No cross-domain, cross-criterion, or model comparison should be interpreted as automatically covered by that adjustment.

The implementation tests exact capacity, score direction, deterministic ties, hand-checkable and tied-score metrics, zero- and all-positive samples, nonfinite scores, clustered resampling, paired comparisons, and exclusion of outcome columns from screening features. Added tests verify prompt separation, training-only scaling, invariance of a held-out fold's predictions to its own labels, single-class fallback, micro-aggregated counts, paired interval reconstruction, and source immutability. Logistic predictions were independently reconstructed with a scikit-learn pipeline. AP and ROC-AUC were checked against scikit-learn 1.8.0. Analysis dependencies include NumPy 2.4.4, pandas 3.0.2, SciPy 1.17.1, and Matplotlib 3.10.8.

The trained models use tolerance $10^{-8}$ and at most 2,000 solver iterations; a convergence warning is an error, not silently accepted. Sorted unique prompt identifiers are permuted with seed 20261002 and divided into five nearly equal prompt-count folds. All policies share stable tie keys within each held-out batch. Bootstrap copies retain the original prompt's fold, with copy-specific deterministic tie keys, preventing the same original prompt from leaking across training and test sets. Alternative partition seeds are 20261003--20261012; no seed is selected based on outcomes. The original-partition models, transformations, fold assignments, out-of-fold scores, and paired bootstrap differences are preserved in the computational record.

\section{Worked Audit-Use Card}\label{app:card}
\begin{description}
\item[Decision and owner.] Proposed evaluator-managed allocation of 132 review slots in a batch of 660 mental-health responses. This is an offline simulation, not an observed institutional workflow.
\item[Measure and rule.] Descending recorded HumT with deterministic seed-based tie resolution. HumT measures human-like tone, not harm probability.
\item[Criterion.] Mean $D\geq3$ under the source rubric; five Group B ratings per response. The trigger includes simulated affect or understanding and does not establish deception experienced by a user.
\item[Coverage.] 42 of 165 criterion cases selected; precision 31.8\%; recall 25.5\% (pointwise 95\% interval 19.2--32.1\%). Random selection catches 33 cases in expectation at the same capacity. The marker comparator catches 63.
\item[Residual concerns.] 123 criterion cases remain among 528 unselected responses. Non-selection is not clearance; responses outside this criterion can contain other concerns.
\item[Uncertainty and limits.] Prompt-cluster bootstrap with 2,000 replicates; dependence on criterion, source population, and a single-turn corpus. No direct evidence of review effectiveness or actual labor savings.
\item[Proposed challenge route.] Permit challenges to the criterion, its application, or allocation separately. Retain disagreements, original ratings, amendments, and reasons. This process remains unevaluated; a low score must not be the sole reason to exclude a concern from consideration.
\end{description}

\section{Reading the Supporting Record}
The computational outputs include matched-budget performance; threshold and rater-consensus sensitivity; criterion intersections and continuous relationships; context counts; complete-prompt model comparisons; join and tie sensitivity; and explicitly hypothetical time, cost, and mixture scenarios. They preserve sparse and undefined cells rather than using favorable subsets for headline claims. New human-study materials exist only as preparation and are not evidence reported in this manuscript.
\end{document}
